% Spira nel piano del foglio, campo esterno che entra e diminuisce (ud6-magnetismo-verso-indotta):
% la corrente indotta circola in senso orario e dentro la spira produce un campo che entra nel foglio.
\begin{tikzpicture}[>=Stealth, font=\small,
    corrente/.style={attenzione, very thick, postaction={decorate},
      decoration={markings, mark=at position 0.125 with {\arrow{Stealth}},
        mark=at position 0.375 with {\arrow{Stealth}}, mark=at position 0.625 with {\arrow{Stealth}},
        mark=at position 0.875 with {\arrow{Stealth}}}}]
  % campo esterno
  % campo esterno, senza i simboli che cadrebbero sul filo della spira
  \foreach \x in {-2.4,-1.6,...,2.5} \foreach \y in {-1.6,-0.8,...,1.7} {
    \pgfmathparse{abs(veclen(\x,\y) - 1.45) > 0.3 ? 1 : 0}
    \ifnum\pgfmathresult=1 \node[medica!55, font=\scriptsize] at (\x,\y) {$\otimes$}; \fi
  }
  % spira, percorsa in senso orario (dall'alto verso destra)
  \draw[corrente] (0,1.45) arc[start angle=90, end angle=-270, radius=1.45];
  \node[attenzione, fill=white, inner sep=1pt] at (0,1.85) {$I$, senso orario};
  % campo della corrente indotta, dentro la spira
  \foreach \p in {(-0.45,0.3),(0.45,0.3),(-0.45,-0.45),(0.45,-0.45)}
    \node[esempio, font=\large] at \p {$\otimes$};
  \node[esempio, fill=white, inner sep=1pt, align=center, font=\footnotesize] at (0,-2.25)
    {campo della corrente indotta: entra nel foglio};
  \node[medica, anchor=west, align=left] at (2.75,0.4) {$\vec B$ esterno:\\entra e diminuisce};
\end{tikzpicture}
